%\thispagestyle{plain}

\chapter*{Abstract}
\addcontentsline{toc}{chapter}{\textit{Abstract}}

%\begin{center}
%\textbf{Abstract}
%\end{center}

Since Dirac first theorized antimatter in 1928 and the discovery of the positron by Anderson in 1932, antimatter has been a subject of study in physics, particularly for investigating fundamental symmetries in nature, such as charge-parity-time reversal symmetry (CPT). After more than 100 years, antimatter is still an area of active interest, and studied both at large scale experiments, such the one operating at the Large Hadron Collider (LHC), and in smaller-scale experiment, such those conducted at the Antimatter-Factory at CERN, including the ALPHA experiment (\textit{Anti-
hydrogen Laser physics Apparatus}).  The ALPHA experiment focuses on producing, confining and conducting spectroscopy as well as gravitational studies on antihydrogen atoms. Those atoms, produced by merging positron and antiproton plasma clouds in a Penning-Malmberg trap, are magnetically confined using superconducting solenoid and octupole magnets. Once a sufficient number of antihydrogen atoms have been accumulated in the trap, spectroscopic experiment are performed. These experiments induce quantum transitions, causing the atoms to change from trappable to untrappable states, which leads to their annihilation on the trap walls. These resulting annihilation events are detected by a silicon vertex detector, and constitutes the experimental spectra, from which the properties of the energy levels of antihydrogen atoms can be deduced. This thesis is dedicated to presenting the analysis and recent results achieved by the ALPHA collaboration, that has measured the hyperfine splitting of the ground state of antihydrogen with a remarkable precision of 4 parts per million (ppm), marking a significant advance in antimatter research. The ground state has been measured in a non-zero magnetic field configuration, where the degeneracy of the 1S energy level if fully resolved by the \SI{1}{\tesla} field. Of the four hyperfine spectral lines, which differ according to the relative spin orientations of the positron and antiproton, only two correspond to trappable states. The hyperfine levels are measured by inducing positron spin-flip transitions from trappable to untrappable states using microwave radiation directed into the trap, at increasing frequency steps. The analysis has consisted on developing the analytic model to interpolate the data, and estimate statistical and systematic uncertainties